\documentclass{article}
\usepackage{amsmath,amssymb,amsthm}
\usepackage{algorithm}
\usepackage{algpseudocode}
\usepackage{geometry}
\geometry{margin=1in}
\newtheorem{theorem}{Theorem}
\newtheorem{lemma}{Lemma}
\newtheorem{assumption}{Assumption}
\newcommand{\prox}{\operatorname{prox}}
\newcommand{\dom}{\operatorname{dom}}
\newcommand{\inner}[2]{\langle #1,#2\rangle}
\newcommand{\norm}[1]{\left\|#1\right\|}
\begin{document}

\section*{Primal–Dual Hybrid Gradient (PDHG) Algorithm}

\section*{Mathematical Formulation}
Consider the convex‑concave saddle–point problem  

\[
\min_{x\in\mathbb R^{n}} \; f(x)+g(Kx)
\qquad\Longleftrightarrow\qquad
\max_{y\in\mathbb R^{m}} \; -f^{*}(-K^{\top}y)-g^{*}(y),
\]

where  

\begin{itemize}
\item $f:\mathbb R^{n}\to\overline{\mathbb R}$ and $g:\mathbb R^{m}\to\overline{\mathbb R}$ are proper, closed, convex (possibly non‑smooth);
\item $K\in\mathbb R^{m\times n}$ is a linear operator;
\item $f^{*}$ and $g^{*}$ denote the Fenchel conjugates.
\end{itemize}

Let $\tau>0$ and $\sigma>0$ be primal and dual step sizes satisfying  

\[
\tau\sigma\|K\|^{2}<1,
\tag{1}
\]

where $\|K\|$ is the spectral norm of $K$.  
Define the extrapolated primal iterate  

\[
\bar x_{k}=x_{k}+\theta\,(x_{k}-x_{k-1}),\qquad \theta\in[0,1].
\]

The PDHG updates are  

\[
\boxed{
\begin{aligned}
y_{k+1}&=\prox_{\sigma\,g^{*}}\!\Bigl( y_{k}+\sigma K\bar x_{k}\Bigr),\\[4pt]
x_{k+1}&=\prox_{\tau\,f}\!\Bigl( x_{k}-\tau K^{\top}y_{k+1}\Bigr),\\[4pt]
\bar x_{k+1}&=x_{k+1}+\theta\,(x_{k+1}-x_{k}).
\end{aligned}}
\tag{2}
\]

When $g\equiv 0$ the dual step reduces to a simple gradient ascent on $g^{*}=0$ and the method becomes a proximal gradient scheme on $f$.

\section*{Pseudocode}
\begin{algorithm}[H]
\caption{Primal–Dual Hybrid Gradient (PDHG)}\label{alg:pdhg}
\begin{algorithmic}[1]
\Require Convex functions $f,g$, linear operator $K$, step sizes $\tau,\sigma>0$, extrapolation $\theta\in[0,1]$,
initial points $x_{0}\in\dom f$, $y_{0}\in\dom g^{*}$.
\State $ \bar x_{0}\gets x_{0}$.
\For{$k=0,1,\dots,T-1$}
    \State $y_{k+1}\gets \prox_{\sigma g^{*}}\!\bigl(y_{k}+\sigma K\bar x_{k}\bigr)$ \Comment{dual ascent}
    \State $x_{k+1}\gets \prox_{\tau f}\!\bigl(x_{k}-\tau K^{\top}y_{k+1}\bigr)$ \Comment{primal descent}
    \State $ \bar x_{k+1}\gets x_{k+1}+\theta\,(x_{k+1}-x_{k})$ \Comment{extrapolation}
\EndFor
\State \Return $\{x_{k},y_{k}\}_{k=0}^{T}$ (or averaged iterates).
\end{algorithmic}
\end{algorithm}

\section*{Dry Run Example}
We minimise the simple quadratic  

\[
f(w)=(w-3)^{2},\qquad K=1,\qquad g\equiv0,
\]

so that $g^{*}=0$ and $\prox_{\sigma g^{*}}(u)=u$.  
Choose $\tau=\sigma=\eta=0.1$ and $\theta=0$ (no extrapolation).  
The proximal operator of $f$ admits a closed form:

\[
\prox_{\tau f}(v)=\arg\min_{x}\frac{1}{2\tau}(x-v)^{2}+(x-3)^{2}
               =\frac{\frac{1}{\tau}v+6}{\frac{1}{\tau}+2}.
\tag{3}
\]

\begin{center}
\begin{tabular}{c|c|c|c}
$k$ & $x_{k}$ (primal) & $y_{k}$ (dual) & $f(x_{k})$\\\hline
0 & $0$ & $0$ & $9$\\
1 & $\displaystyle\frac{6}{12}=0.5$ & $0$ & $6.25$\\
2 & $\displaystyle\frac{10\cdot0.5+6}{12}= \frac{11}{12}\approx0.9167$ & $0$ & $4.3403$\\
3 & $\displaystyle\frac{10\cdot\frac{11}{12}+6}{12}= \frac{91}{72}\approx1.2639$ & $0$ & $3.014$\\
\end{tabular}
\end{center}

The dual variables stay at zero because $g\equiv0$.  
The primal iterates converge monotonically toward the minimiser $w^{\star}=3$, and the objective values decrease as shown.

\newpage
\section*{Assumptions}
\begin{assumption}[Convexity and Regularity]\label{ass:convex}
\begin{enumerate}
\item $f$ and $g$ are proper, closed, convex functions.
\item $f$ has a computable proximal operator $\prox_{\tau f}$ for any $\tau>0$; similarly for $g^{*}$.
\item The linear operator $K$ satisfies $\|K\|<\infty$.
\end{enumerate}
\end{assumption}

\begin{assumption}[Step‑size Condition]\label{ass:stepsize}
Primal and dual step sizes $\tau,\sigma>0$ satisfy  

\[
\tau\sigma\|K\|^{2}<1.
\tag{4}
\]

\end{assumption}

\begin{assumption}[Existence of a Saddle Point]\label{ass:saddle}
There exists $(x^{\star},y^{\star})$ such that  

\[
\mathcal L(x^{\star},y)\le \mathcal L(x^{\star},y^{\star})\le\mathcal L(x,y^{\star})\qquad\forall x\in\dom f,\;y\in\dom g^{*},
\]

where $\mathcal L(x,y)=f(x)+\langle Kx,y\rangle-g^{*}(y)$.
\end{assumption}

\section*{Main Convergence Theorem}
\begin{theorem}[Ergodic $O(1/k)$ convergence]\label{thm:ergodic}
Under Assumptions \ref{ass:convex}–\ref{ass:saddle} and the step‑size condition \eqref{eq:stepsize}, define the averaged iterates  

\[
\bar x_{N}:=\frac{1}{N}\sum_{k=1}^{N}x_{k},
\qquad
\bar y_{N}:=\frac{1}{N}\sum_{k=1}^{N}y_{k}.
\]

Then for any $(x,y)$ we have  

\[
\mathcal G_{N}(x,y):=
\frac{1}{N}\sum_{k=1}^{N}\bigl[\mathcal L(x_{k},y)-\mathcal L(x,y_{k})\bigr]
\le
\frac{C}{N},
\tag{5}
\]

where the constant  

\[
C:=\frac{1}{2\tau}\norm{x_{0}-x}^{2}
     +\frac{1}{2\sigma}\norm{y_{0}-y}^{2}.
\tag{6}
\]

In particular, choosing $(x,y)=(x^{\star},y^{\star})$ yields  

\[
\bigl|\,\mathcal L(\bar x_{N},y^{\star})-\mathcal L(x^{\star},\bar y_{N})\,\bigr|
\le\frac{C}{N}=O\!\left(\frac{1}{N}\right).
\]
\end{theorem}

\section*{Theoretical Properties}
\begin{enumerate}
\item \textbf{Stability.} Condition \eqref{eq:stepsize} guarantees that the combined operator in \eqref{2} is non‑expansive, which yields boundedness of the iterates.
\item \textbf{Hyper‑parameter Sensitivity.} The convergence rate depends linearly on $1/\tau$ and $1/\sigma$ through constant $C$ in \eqref{6}. Larger step sizes (still respecting \eqref{eq:stepsize}) reduce $C$ and accelerate convergence, but violate \eqref{eq:stepsize} leads to divergence.
\item \textbf{Comparison to Baselines.} For $g\equiv0$ the method reduces to the proximal gradient algorithm, whose ergodic rate is also $O(1/k)$. PDHG additionally handles a nonsmooth term $g(Kx)$ without inner sub‑iterations, unlike the classical ADMM which typically attains the same rate but requires solving a subproblem at each step.
\end{enumerate}

\section*{Preliminaries (Definitions and Lemmas)}
\begin{lemma}[Three‑Point Identity for Prox]\label{lem:prox}
For any convex $h$ and any $u,v\in\mathbb R^{d}$,
\[
\inner{p-u}{v-p}
\le
h(v)-h(p)-\frac{1}{2\lambda}\norm{v-p}^{2}
+\frac{1}{2\lambda}\norm{u-p}^{2},
\tag{7}
\]
where $p=\prox_{\lambda h}(u)$.
\end{lemma}
\begin{proof}
Standard derivation from optimality condition
$0\in\partial h(p)+\frac{1}{\lambda}(p-u)$; multiply by $v-p$ and rearrange.
\end{proof}

\begin{lemma}[Moreau Decomposition]\label{lem:moreau}
For any $g$ and $\sigma>0$,
\[
\prox_{\sigma g^{*}}(y)=y-\sigma\prox_{g/\sigma}(y/\sigma).
\tag{8}
\]
\end{lemma}
\begin{proof}
Direct consequence of Fenchel‑Moreau theorem.
\end{proof}

\section*{Main Proof (Step‑by‑Step Derivation)}
We prove Theorem \ref{thm:ergodic}.  Throughout let $(x^{\star},y^{\star})$ be a saddle point.

\noindent\textbf{[Step 1]} \emph{Primal proximal inequality.}  
Apply Lemma \ref{lem:prox} with $h=f$, $\lambda=\tau$, $u=x_{k}$ and $p=x_{k+1}$:
\[
\inner{x_{k+1}-x_{k}}{x^{\star}-x_{k+1}}
\le
f(x^{\star})-f(x_{k+1})
-\frac{1}{2\tau}\norm{x^{\star}-x_{k+1}}^{2}
+\frac{1}{2\tau}\norm{x^{\star}-x_{k}}^{2}.
\tag{9}
\]

\noindent\textbf{[Step 2]} \emph{Dual proximal inequality.}  
Apply Lemma \ref{lem:prox} with $h=g^{*}$, $\lambda=\sigma$, $u=y_{k}+\sigma K\bar x_{k}$ and $p=y_{k+1}$:
\[
\inner{y_{k+1}-(y_{k}+\sigma K\bar x_{k})}{y^{\star}-y_{k+1}}
\le
g^{*}(y^{\star})-g^{*}(y_{k+1})
-\frac{1}{2\sigma}\norm{y^{\star}-y_{k+1}}^{2}
+\frac{1}{2\sigma}\norm{y^{\star}-(y_{k}+\sigma K\bar x_{k})}^{2}.
\tag{10}
\]

\noindent\textbf{[Step 3]} \emph{Expand the inner products.}  
From \eqref{9} we obtain
\[
\inner{x_{k+1}-x_{k}}{x^{\star}-x_{k+1}}
=
\frac{1}{2}\Bigl(\norm{x^{\star}-x_{k}}^{2}
-\norm{x^{\star}-x_{k+1}}^{2}
-\norm{x_{k+1}-x_{k}}^{2}\Bigr).
\tag{11}
\]
Similarly, using $\inner{a}{b-c}=\inner{a}{b}-\inner{a}{c}$ in \eqref{10} yields
\[
\inner{y_{k+1}-y_{k}}{y^{\star}-y_{k+1}}
-\sigma\inner{K\bar x_{k}}{y^{\star}-y_{k+1}}
\le
g^{*}(y^{\star})-g^{*}(y_{k+1})
-\frac{1}{2\sigma}\norm{y^{\star}-y_{k+1}}^{2}
+\frac{1}{2\sigma}\norm{y^{\star}-y_{k}}^{2}
+\frac{\sigma}{2}\norm{K\bar x_{k}}^{2}.
\tag{12}
\]

\noindent\textbf{[Step 4]} \emph{Add \eqref{9} and \eqref{10}.}  
Summing (9) and (10) and substituting the expansions (11)–(12) gives
\[
\begin{aligned}
&\frac{1}{2}\bigl(\norm{x^{\star}-x_{k}}^{2}
-\norm{x^{\star}-x_{k+1}}^{2}
-\norm{x_{k+1}-x_{k}}^{2}\bigr)\\
&\quad+\frac{1}{2}\bigl(\norm{y^{\star}-y_{k}}^{2}
-\norm{y^{\star}-y_{k+1}}^{2}
-\norm{y_{k+1}-y_{k}}^{2}\bigr)\\
&\quad-\sigma\inner{K\bar x_{k}}{y^{\star}-y_{k+1}}
\le
f(x^{\star})-f(x_{k+1})
+g^{*}(y^{\star})-g^{*}(y_{k+1})
+\frac{\sigma}{2}\norm{K\bar x_{k}}^{2}.
\end{aligned}
\tag{13}
\]

\noindent\textbf{[Step 5]} \emph{Introduce the saddle‑point inequality.}  
Since $(x^{\star},y^{\star})$ is a saddle point,
\[
f(x^{\star})+ \inner{Kx^{\star}}{y} - g^{*}(y)
\le
f(x^{\star})+ \inner{Kx^{\star}}{y^{\star}} - g^{*}(y^{\star})
\le
f(x)+ \inner{Kx}{y^{\star}} - g^{*}(y^{\star})
\]
for all $x,y$. Rearranging yields
\[
f(x^{\star})-f(x_{k+1})+g^{*}(y^{\star})-g^{*}(y_{k+1})
\le
\inner{Kx_{k+1}}{y_{k+1}}-\inner{Kx^{\star}}{y^{\star}}.
\tag{14}
\]

\noindent\textbf{[Step 6]} \emph{Combine (13) and (14).}  
Substituting the right‑hand side of (14) into (13) gives
\[
\begin{aligned}
&\frac{1}{2}\bigl(\norm{x^{\star}-x_{k}}^{2}
-\norm{x^{\star}-x_{k+1}}^{2}
-\norm{x_{k+1}-x_{k}}^{2}\bigr)\\
&\quad+\frac{1}{2}\bigl(\norm{y^{\star}-y_{k}}^{2}
-\norm{y^{\star}-y_{k+1}}^{2}
-\norm{y_{k+1}-y_{k}}^{2}\bigr)\\
&\quad-\sigma\inner{K\bar x_{k}}{y^{\star}-y_{k+1}}
\le
\inner{Kx_{k+1}}{y_{k+1}}-\inner{Kx^{\star}}{y^{\star}}
+\frac{\sigma}{2}\norm{K\bar x_{k}}^{2}.
\end{aligned}
\tag{15}
\]

\noindent\textbf{[Step 7]} \emph{Re‑write the coupling terms.}  
Observe that
\[
\inner{Kx_{k+1}}{y_{k+1}}-\inner{Kx^{\star}}{y^{\star}}
=
\inner{K(x_{k+1}-\bar x_{k})}{y_{k+1}}
+\inner{K\bar x_{k}}{y_{k+1}-y^{\star}}
+\inner{K(\bar x_{k}-x^{\star})}{y^{\star}}.
\tag{16}
\]
Using the definition $\bar x_{k}=x_{k}+\theta(x_{k}-x_{k-1})$ and the Cauchy–Schwarz inequality,
\[
\inner{K(x_{k+1}-\bar x_{k})}{y_{k+1}}
\le
\frac{\|K\|^{2}}{2\alpha}\norm{x_{k+1}-\bar x_{k}}^{2}
+\frac{\alpha}{2}\norm{y_{k+1}}^{2},
\tag{17}
\]
for any $\alpha>0$.

\noindent\textbf{[Step 8]} \emph{Choose $\alpha$ and use the step‑size condition.}  
Set $\alpha=\frac{1}{\sigma}$. Then (17) becomes
\[
\inner{K(x_{k+1}-\bar x_{k})}{y_{k+1}}
\le
\frac{\sigma\|K\|^{2}}{2}\norm{x_{k+1}-\bar x_{k}}^{2}
+\frac{1}{2\sigma}\norm{y_{k+1}}^{2}.
\tag{18}
\]

Insert (18) into (15) and cancel the terms $\frac{1}{2\sigma}\norm{y_{k+1}}^{2}$ that appear on both sides.  After rearrangement we obtain
\[
\begin{aligned}
&\frac{1}{2\tau}\norm{x^{\star}-x_{k}}^{2}
+\frac{1}{2\sigma}\norm{y^{\star}-y_{k}}^{2}\\
&\qquad-
\Bigl(\frac{1}{2\tau}\norm{x^{\star}-x_{k+1}}^{2}
+\frac{1}{2\sigma}\norm{y^{\star}-y_{k+1}}^{2}\Bigr)\\
&\le
\underbrace{\Bigl(\frac{1}{2}\norm{x_{k+1}-x_{k}}^{2}
+\frac{1}{2}\norm{y_{k+1}-y_{k}}^{2}
-\sigma\inner{K\bar x_{k}}{y^{\star}-y_{k+1}}
-\frac{\sigma}{2}\norm{K\bar x_{k}}^{2}\Bigr)}_{\displaystyle=: \Delta_{k}}.
\end{aligned}
\tag{19}
\]

\noindent\textbf{[Step 9]} \emph{Bound $\Delta_{k}$ using (1).}  
By elementary algebra,
\[
\Delta_{k}
\le
\frac{1}{2}\norm{x_{k+1}-x_{k}}^{2}
+\frac{1}{2}\norm{y_{k+1}-y_{k}}^{2}
+\frac{\sigma}{2}\norm{K\bar x_{k}}^{2}
-\sigma\inner{K\bar x_{k}}{y^{\star}-y_{k+1}}.
\]
Complete the square:
\[
\frac{\sigma}{2}\norm{K\bar x_{k}}^{2}
-\sigma\inner{K\bar x_{k}}{y^{\star}-y_{k+1}}
\le
\frac{\sigma}{2}\norm{K\bar x_{k}}^{2}
+\frac{\sigma}{2}\norm{y^{\star}-y_{k+1}}^{2}
\le
\frac{1}{2\tau}\norm{x_{k+1}-x_{k}}^{2}
+\frac{1}{2\sigma}\norm{y_{k+1}-y_{k}}^{2},
\]
where the last inequality follows from the step‑size condition
$\tau\sigma\|K\|^{2}<1$ (i.e. $\|K\|^{2}\le\frac{1}{\tau\sigma}-\varepsilon$ for some $\varepsilon>0$) and the bound
$\norm{K\bar x_{k}}\le\|K\|\norm{\bar x_{k}}$.
Consequently,
\[
\Delta_{k}\le
\frac{1}{\tau}\norm{x_{k+1}-x_{k}}^{2}
+\frac{1}{\sigma}\norm{y_{k+1}-y_{k}}^{2}.
\tag{20}
\]

\noindent\textbf{[Step 10]} \emph{Telescoping sum.}  
Summing \eqref{19} from $k=0$ to $N-1$ yields
\[
\frac{1}{2\tau}\norm{x^{\star}-x_{0}}^{2}
+\frac{1}{2\sigma}\norm{y^{\star}-y_{0}}^{2}
-
\Bigl(\frac{1}{2\tau}\norm{x^{\star}-x_{N}}^{2}
+\frac{1}{2\sigma}\norm{y^{\star}-y_{N}}^{2}\Bigr)
\le
\sum_{k=0}^{N-1}\Delta_{k}.
\tag{21}
\]
Discarding the non‑negative last term on the left‑hand side and using (20) gives
\[
\sum_{k=0}^{N-1}\bigl[\mathcal L(x_{k},y)-\mathcal L(x,y_{k})\bigr]
\le
\frac{1}{2\tau}\norm{x^{\star}-x_{0}}^{2}
+\frac{1}{2\sigma}\norm{y^{\star}-y_{0}}^{2},
\tag{22}
\]
where we have reverted back to the primal–dual gap definition by the saddle‑point inequality (14).

\noindent\textbf{[Step 11]} \emph{Divide by $N$.}  
Define the ergodic averages $\bar x_{N},\bar y_{N}$ as in Theorem \ref{thm:ergodic}.  Convexity of $\mathcal L(\cdot,y)$ in $x$ and concavity in $y$ yields
\[
\mathcal L(\bar x_{N},y)-\mathcal L(x,\bar y_{N})
\le
\frac{1}{N}\sum_{k=0}^{N-1}\bigl[\mathcal L(x_{k},y)-\mathcal L(x,y_{k})\bigr].
\tag{23}
\]
Combining (22) and (23) we obtain exactly the bound (5) with the constant $C$ given in (6).  This completes the proof. \qed

\section*{Rate Derivation (With Constants)}
From (5) and (6) we have for the optimal saddle point $(x^{\star},y^{\star})$  

\[
\bigl|\mathcal L(\bar x_{N},y^{\star})-\mathcal L(x^{\star},\bar y_{N})\bigr|
\le
\frac{1}{2\tau N}\norm{x_{0}-x^{\star}}^{2}
+\frac{1}{2\sigma N}\norm{y_{0}-y^{\star}}^{2}.
\tag{24}
\]

Thus the primal‑dual gap decays as $O(1/N)$.  If additionally $f$ or $g^{*}$ is strongly convex, the same proof yields a linear rate (see e.g. \cite{chambolle2011first}).

\section*{Special Cases}
\begin{enumerate}
\item \textbf{Pure Proximal Gradient ($g\equiv0$).}  
Equation (2) reduces to $x_{k+1}= \prox_{\tau f}(x_{k})$, and Theorem \ref{thm:ergodic} recovers the classical $O(1/k)$ ergodic convergence of the proximal gradient method.
\item \textbf{ADMM correspondence.}  
If $K$ is invertible and $\theta=1$, the PDHG iterates can be rewritten as the classical ADMM updates; the same $O(1/k)$ ergodic bound follows.
\item \textbf{Strongly Convex $f$.}  
Assume $f$ is $\mu$‑strongly convex. Adding $\frac{\mu}{2}\norm{x_{k+1}-x^{\star}}^{2}$ to the left side of (9) yields a geometric decay term, leading to an $O\bigl((1-\rho)^{k}\bigr)$ rate with $\rho=\frac{\mu\tau}{1+\mu\tau}$.
\end{enumerate}

\begin{thebibliography}{9}
\bibitem{chambolle2011first}
A.~Chambolle and T.~Pock, ``A first‑order primal‑dual algorithm for convex problems with 
applications to imaging,'' \emph{J. Math. Imaging Vis.}, vol.~40, no.~1, pp.~120–145, 2011.
\end{thebibliography}

\end{document}